% ============================================================================================================
% Abraham's contribution, part B: the "Technical Implementation" section, revised.
% Replaces the draft's subsections from "Cross-Validation" through "Model Evaluation" (keep the draft's
% "Prediction Problem", "Constructed Pair-Level Predictors", "Pre-Date Feature Sets" and "Survey-Format
% Standardization" as they are, with the small fixes listed in report/README.md). Every implementation claim was
% checked against the code (report/factcheck_draft.md); citation keys are in report/references_abraham.bib.
% ============================================================================================================

\subsection{Cross-Validation}

A central feature of the analysis is that participants appear in multiple dyads. Randomly splitting dyads into
training and validation sets would place the same participant on both sides of the split, so the validation error
would measure how well a model predicts people it has already seen. We therefore used grouped cross-validation
(Section~\ref{sec:beyond}): five folds formed from whole experimental waves. The 21 events were ordered at random
with a fixed seed (20261006) and dealt into five folds of five, four, four, four and four events. In each iteration
every dyad of the held-out events was set aside and the model was estimated on the remaining events, so no
participant in a held-out event appeared in that iteration's training data. The held-out predictions of the five
folds were then pooled to compute each performance measure, with each dyad classified at its own fold's threshold.

For models with a tuning parameter we used nested cross-validation \cite{Varma:2006}. Within each outer training
set, the training events were divided into four further event-grouped folds; each candidate value was scored by the
pooled AUC over these inner folds, with all preprocessing refitted inside each inner training set, and the best value
was used to refit the model on the whole outer training set before it scored the outer held-out events. The outer
held-out events are therefore never used to choose a tuning parameter.

\subsection{Preprocessing Within Cross-Validation}

Every preprocessing step that learns from the data was fitted on the training fold only and applied unchanged to
the held-out fold, which is the ``right way'' to cross-validate a pipeline \cite[Section 7.10.2, pp.~245--247]{Hastie:2009}.

\begin{enumerate}
\item \emph{Skewed predictors.} A numeric predictor with more than two distinct values was transformed to
$\log(1+x-\min x)$ if its training-fold skewness exceeded $1$, and to $-\log(1+\max x-x)$ if it was below $-1$. For
person-level variables the skewness was computed over one value per participant. The minimum and maximum came from
the training fold, and held-out values of transformed predictors were clipped to that range.
\item \emph{Missing values.} Missing numeric and binary values were filled with the training-fold median. Field and
career codes held by fewer than 10 participants in the training fold were merged into ``other'', and categorical
variables were one-hot coded with the most common training level as the reference.
\item \emph{Standardization.} Numeric columns, including any spline columns, were standardized with the training
mean and standard deviation; dummy variables and binary indicators were left as $0/1$.
\end{enumerate}

The survey-format standardization of the stated preferences (described above) is the one exception: it was computed
once, before cross-validation, from all participants' survey answers. It uses no outcomes, so it cannot leak match
information into a fold.

\subsection{Prediction Models}

\subsubsection{Logistic regression, with optional splines}
Our primary model was multiple logistic regression \cite[Equation 4.7, p.~142]{James:2023},
\[
P(Y_i=1\mid X_i)=\frac{1}{1+\exp\left[-(\beta_0+X_i^\top\beta)\right]},
\]
fitted by maximum likelihood with a very small ridge penalty ($\lambda=10^{-6}$, intercept unpenalized) to keep the
estimates finite when some predictor combinations nearly separate matches from non-matches. Six continuous
predictors (the two ages, the absolute age difference, interest correlation and the two preference-fit measures)
could each enter as a natural cubic spline with three degrees of freedom instead of a linear term
\cite[Section 7.4.3, p.~297]{James:2023}, \cite[Section 5.2.1, pp.~144--146]{Hastie:2009},
if, tried one at a time, the spline lowered the training-fold Bayesian information criterion
\[
\mathrm{BIC}=-2\ell+k\log n
\]
\cite[Section 7.7, Eq.~(7.35), p.~233]{Hastie:2009}, the likelihood form of the criterion in
\cite[Section 6.1.3, Equation 6.3, p.~237]{James:2023},
where $\ell$ is the maximized log-likelihood, $k$ the number of parameters and $n$ the number of training dyads. No spline
was selected in any fold, so the reported full logistic model is linear.

\subsubsection{Stepwise logistic regression}
We also selected predictors by BIC with forward, backward and both-directions stepwise search
\cite[Section 6.1.2, pp.~233--235]{James:2023}.
Selection worked on whole predictors: the dummy variables of a categorical predictor, and the columns of a spline,
entered or left together. Forward and both-directions search started from the intercept-only model, backward search
from the full model, and each stopped when no single addition or removal lowered BIC.

\subsubsection{Linear and quadratic discriminant analysis}
LDA used the standard estimates of Section~\ref{sec:lda_theory}, with class priors equal to the training-fold match
and non-match proportions. For QDA, estimating a separate $124\times124$ covariance matrix for the roughly 550 training
matches is unstable, so we used regularized QDA \cite{Friedman:1989}, shrinking each class covariance toward the
identity \cite[Section 4.3.1, Eq.~(4.14), p.~112]{Hastie:2009},
\[
\widetilde\Sigma_k=(1-r)\hat\Sigma_k+rI,\qquad r\in\{0.1,0.3,0.5,0.7,0.9\},
\]
with $r$ chosen by the inner event-grouped cross-validation. The strongest shrinkage, $r=0.9$, was chosen in every
fold, the edge of the grid.

\subsubsection{Random forest}
The random forest (Section~\ref{sec:random_forest_theory}; \cite{Breiman:2001}, \cite[Algorithm 15.1, p.~588]{Hastie:2009})
used 500 classification trees,
$\sqrt p$ predictors considered at each split and the Gini criterion, and predicted a match probability by averaging
the trees. The minimum leaf size was tuned over $\{1,5,20,50\}$ by inner cross-validation; on the pre-date predictors
the largest value, 50, was chosen in every fold.

\subsubsection{L1-regularized logistic regression}
The L1 (lasso) penalty \cite[Section 6.2.2, pp.~244--252]{James:2023} shrinks some coefficients exactly to zero; for
logistic regression it is added to the negative log-likelihood \cite[Section 4.4.4, Eq.~(4.31), p.~125]{Hastie:2009}.
We parametrized the penalty by its inverse strength $C$ and searched 30 log-spaced values from
\[
C_{\min}=\frac{1}{\max_j\left|\sum_i D_{ij}(Y_i-\bar Y)\right|},
\]
the value at which every coefficient is zero (the logistic analogue of \cite[Section 2.5, p.~7]{Friedman:2010}),
to $1000\,C_{\min}$, where $D$ is the standardized training design matrix. The value with the largest pooled inner-fold
AUC was selected. This model was added after the pre-registered comparisons and is reported as an additional
specification.

\subsection{Model Evaluation}

Because only 16.5\% of dyads resulted in a match, a classifier that predicts ``no match'' for everyone is 83.5\%
accurate, so we report sensitivity, specificity, balanced accuracy and, as the principal measure, the AUC. For the
threshold-based measures a dyad is a predicted match when its predicted probability exceeds the match rate in its
training fold (0.158 to 0.171 across folds).

Because participants appear in several dyads and all dyads of an event share the same evening, observations from the
same event are not independent. Every standard error in this report is therefore clustered by event: the
cluster-robust variance for error rates, sensitivity, specificity and balanced accuracy, and Obuchowski's clustered
variance for the AUC (Section~\ref{sec:beyond}). We checked these standard errors by simulation before using them:
with 21 events the AUC standard errors were 0.84 to 0.95 of the true spread, so a nominal 90\% interval behaves like
roughly an 83\% interval.
